\documentclass[11pt]{article}
\usepackage[a4paper,margin=27mm]{geometry}
\usepackage{amsmath,amssymb,amsthm,hyperref}
\setlength{\emergencystretch}{2em}
\newtheorem{theorem}{Theorem}
\newtheorem{lemma}[theorem]{Lemma}
\newtheorem{corollary}[theorem]{Corollary}
\newcommand{\Q}{\mathbb Q}\newcommand{\Z}{\mathbb Z}
\title{An obstruction to compressing tuple integrality by polynomial tests}
\author{\(\mathbb Z\mathbb Q\) Diophantine Project}
\date{23 November 2026}
\begin{document}\maketitle
\begin{abstract}
Fewer than $r$ polynomial integer-value tests cannot recognize all integral
points of $\Q^r$, even when rational coefficients are allowed. The
polynomials need not be homogeneous. In fact, for
infinitely many rational primes there are nonintegral tuples, with no
other denominator primes, on which all the tests are integral. The proof
uses a positive-dimensional algebraic fiber, a branch at infinity, and
$p$-adic approximation. This rules out a broad class of direct certificate
compressions; it says nothing against sharing existential witnesses or
using additional equations in a universal system.
\end{abstract}

\begin{theorem}\label{main}
Let $0\le k<r$, and let $F_1,\ldots,F_k\in\Z[X_1,\ldots,X_r]$.
There are infinitely many rational primes $p$ for which there are tuples
\[
 c\in\Z[1/p]^r\setminus\Z^r,
 \qquad F_i(c)\in\Z\quad(1\le i\le k).
\]
For each of these primes, the denominator exponents of such tuples are
unbounded. Consequently an equivalence
\[
 x\in\Z^r\quad\Longleftrightarrow\quad
 F_1(x),\ldots,F_k(x)\in\Z
 \qquad(x\in\Q^r)
\]
requires $k\ge r$.
\end{theorem}

We first give the geometric and local details used in the proof.
The elementary case $k=0$ is also covered by the argument.

\begin{lemma}[Infinitely many degree-one places]\label{places}
For every number field $K$ and every finite set $S$ of rational primes,
there is a rational prime $p\notin S$ and an embedding $K\hookrightarrow\Q_p$.
The primes admitting such an embedding form an infinite set.
\end{lemma}
\begin{proof}
Choose an integral primitive element for $K$, with monic minimal polynomial
$f\in\Z[T]$. A nonconstant integral polynomial has infinitely many prime
divisors of its nonzero integer values. Here is Schur's elementary argument.
If the set of such primes were finite, translate the variable to make
$a=f(0)\ne0$, and let $M$ be their product. For an integer $t$,
\[
 f(aMt)=a\bigl(1+Mt\,b(t)\bigr)
\]
with $b\in\Z[T]$. The second factor is congruent to one at every alleged
prime divisor. For some $t$ its absolute value exceeds one, so it has a
prime divisor outside that set, a contradiction.

Apply this fact after excluding $S$ and the finitely many divisors of the
discriminant of $f$. For any remaining prime divisor $p$ of $f(n)$, the
residue of $n$ is a simple root of $f$ modulo $p$. Hensel lifting supplies
a root in $\Z_p$. Evaluation at this root gives an embedding
$K\simeq\Q[T]/(f)\hookrightarrow\Q_p$. The same proof works after
excluding any further finite set.
\end{proof}

\begin{lemma}[An unbounded local branch]\label{branch}
Let $V\subset\mathbb A^r_{\Q}$ be an affine algebraic set having a
positive-dimensional geometric component. There are a number field $K$,
a finite set $S$ of rational primes, a smooth projective curve $C/K$,
a point $P\in C(K)$, and rational functions $x_1,\ldots,x_r\in K(C)$
with the following properties.
\begin{enumerate}
\item On a nonempty open subset of $C$, the tuple $(x_1,\ldots,x_r)$ lies
in $V$.
\item Some $x_j$ has a pole at $P$.
\item For every embedding $K\hookrightarrow\Q_p$ with $p\notin S$,
there are $\Q_p$-points in that open subset, arbitrarily close to $P$,
on which $v_p(x_j)$ tends to $-\infty$.
\end{enumerate}
\end{lemma}
\begin{proof}
Choose a geometrically irreducible affine curve on a positive-dimensional
geometric component,
by intersecting with sufficiently many general hyperplanes. All the
finitely many coefficients involved are algebraic, so the curve is
defined over a number field. Take the normalization of its projective
closure and enlarge the field if necessary. In characteristic zero this
is a smooth projective curve $C$. The original affine coordinate functions
are rational functions $x_i$ on $C$.

At least one of these functions is nonconstant. A nonconstant rational
function on a projective curve has a pole; otherwise it would be a global
regular function and hence constant. Choose such a pole $P$, and enlarge
$K$ once more so that $P$ is $K$-rational. The removed points of the
original affine curve form a finite set on $C$.

Choose a local parameter $t\in K(C)$ at $P$. If $e=-\operatorname{ord}_P(x_j)>0$,
then locally
\[
                         x_j=t^{-e}u,
\]
where $u$ is regular and $u(P)\ne0$. Shrink an affine neighborhood of $P$
so that $u$ is a unit and $t$ is an etale coordinate there. These choices
involve only finitely many coefficients and nonzero determinants. Exclude
the rational primes dividing the norms of their denominators and nonzero
constant values, as well as the finitely many primes of bad reduction. This is a finite set $S$.

After any embedding at a remaining prime, the implicit function form of
Hensel's lemma gives a $p$-adic neighborhood of $P$ parameterized by an
open ball around $t=0$. In particular it contains points with $t=p^n$
for every sufficiently large $n$. The value $u$ is a $p$-adic unit on a
sufficiently small such neighborhood, so
\[
                         v_p(x_j)=-en.
\]
Nonzero sufficiently small $t$ avoids the finitely many removed points.
This proves all three assertions. The standard facts used here are
normalization of a projective curve, the existence of a uniformizer at a
smooth rational point, and Hensel's local inverse theorem; no rational
parameterization of $C$ is assumed.
\end{proof}

\begin{proof}[Proof of Theorem~\ref{main}]
Consider the fiber through the origin
\[
 V=\{x:F_i(x)=F_i(0)\text{ for }1\le i\le k\}.
\]
It is nonempty. The height bound for an ideal generated by $k$ elements
in $\overline{\Q}[X_1,\ldots,X_r]$ gives every irreducible component
dimension at least $r-k>0$. Apply Lemma~\ref{branch}. By
Lemma~\ref{places}, infinitely many primes admit the required embeddings
and avoid its finite exceptional set.

Fix one such prime. The branch provides a tuple $x\in V(\Q_p)$ with a
coordinate of arbitrarily negative valuation. Since $F_i(x)=F_i(0)\in\Z$,
continuity gives a $p$-adic neighborhood of $x$ on which every $F_i$ takes
values in $\Z_p$. Shrink this neighborhood so that the chosen coordinate
has exactly the same negative valuation as at $x$.

The ring $\Z[1/p]$ is dense in $\Q_p$: finite truncations of a $p$-adic
expansion lie in this ring. Hence this neighborhood contains
$c\in\Z[1/p]^r$. Each $F_i(c)$ belongs to $\Z[1/p]$, because $F_i$ has
integer coefficients, and also belongs to $\Z_p$. The elementary identity
\[
                         \Z[1/p]\cap\Z_p=\Z
\]
therefore gives $F_i(c)\in\Z$. The retained negative valuation proves
$c\notin\Z^r$. Its possible denominator exponent tends to infinity along
the branch. This proves the theorem.
\end{proof}

\begin{corollary}[Rational-coefficient tests]\label{rational-coefficients}
Let $0\le k<r$ and $F_1,\ldots,F_k\in\Q[X_1,\ldots,X_r]$.
Suppose there is an integral tuple $a\in\Z^r$ such that
$F_i(a)\in\Z$ for every $i$. Then the conclusion of
Theorem~\ref{main} holds: for infinitely many rational primes $p$ there
are nonintegral tuples $c\in\Z[1/p]^r$ with all $F_i(c)\in\Z$, and
their $p$-denominator exponents are unbounded.

Consequently no family of fewer than $r$ rational-coefficient
polynomials recognizes $\Z^r$ by simultaneous integrality of its values.
This includes integer-valued polynomials, whose coefficients need not
themselves be integers.
\end{corollary}
\begin{proof}
Choose a positive integer $D$ clearing the denominators of all
nonconstant coefficients of all translated polynomials $F_i(a+X)$.
Their constant terms are already integers. Hence
\[
                       G_i(X)=F_i(a+DX)\in\Z[X_1,\ldots,X_r],
\]
because each nonconstant monomial acquires at least one factor of $D$.
Apply Theorem~\ref{main} to $G_1,\ldots,G_k$ and discard the finitely
many primes dividing $D$. For every remaining prime and every resulting
tuple $b$, set $c=a+Db$. Then $c\in\Z[1/p]^r$ and
$F_i(c)=G_i(b)\in\Z$. Whenever $v_p(b_j)<0$, integrality of $a_j$ and
$v_p(D)=0$ give $v_p(c_j)=v_p(b_j)$. Thus nonintegrality and arbitrarily
large denominator exponents are preserved. A purported recognition of
$\Z^r$ automatically satisfies the hypothesis at $a=0$, which proves
the final assertion.
\end{proof}

\section*{Exact scope for the undecidability project}
The theorem strengthens the homogeneous obstruction in the general-degree
transfer note: it permits arbitrary degrees and lower-order terms;
Corollary~\ref{rational-coefficients} also permits rational coefficients.
It
shows that one cannot repair the failure merely by excluding finitely
many denominator primes. It is a lower bound only for the stated model
of direct polynomial integer-value tests on all of $\Q^r$.

It does not forbid simultaneous integer-defining systems with shared
elliptic points, contact parameters, or finite-algebra encoders. Those
systems have additional existential witnesses. It also does not apply
automatically to a tuple constrained to solve a universal quadratic
system: the final rational approximation in the proof preserves
integrality of the polynomial values but need not preserve other
polynomial equalities exactly. No global lower bound for the number of
variables in rational undecidability follows.
\end{document}
